% Generated by report/make_restated.py from theory/B3_prop59_tightness.tex: do not edit (edit the note and rerun).
\begin{restated}{Corollary}{cor:B3-c}{B3}{1}{parametrization $\Delta=c\gamma^2$}
Write $\Delta=c\gamma^2$ with $c>2$. Then $\alpha=1-\gamma R^2(1+\frac c2)$, the condition $\alpha>0$ reads $\gamma<\frac{2}{R^2(2+c)}$, and \eqref{ineq:B3-general} becomes
\begin{equation*}\tag{\ref*{ineq:B3-general-c}}
\begin{aligned}
\Ei[V_t]&\le\big(1-\min\{\tau_D(\gamma,c),\tau_M(c)\}\big)V_{t-1},\\
\tau_D(\gamma,c)&:=2\gamma\mu\Big(1-\gamma R^2\Big(1+\frac c2\Big)\Big),\qquad
\tau_M(c):=\frac{1-2/c}{n}.
\end{aligned}
\end{equation*}
For fixed $c$, $\gamma\mapsto\tau_D(\gamma,c)$ is a concave parabola, maximal at
\[
\gamma^\star(c)=\frac{1}{R^2(2+c)},\qquad \tau_D(\gamma^\star(c),c)=\frac{\mu}{R^2(2+c)}=\frac{1}{\kappa(2+c)} ,
\]
while $\tau_M$ does not depend on $\gamma$. Hence the best rate that \eqref{ineq:B3-general-c} certifies for a given $c$ is
\postdisplaypenalty=10000 \[
\tau(c):=\min\Big\{\frac1{\kappa(2+c)},\ \frac{c-2}{cn}\Big\}.
\]
\end{restated}
